\documentclass[a4paper,12pt]{article}

% 字体与页面
\usepackage[UTF8]{ctex}
\usepackage{fontspec}
\setmainfont{Times New Roman}

\usepackage{geometry}
\geometry{left=2.5cm, right=2.5cm, top=2.5cm, bottom=2.5cm}
\linespread{1.2} % n 倍行距
\setlength{\parskip}{0em} % 段间距紧凑

% 页眉页脚
\usepackage{fancyhdr}
\pagestyle{fancy}
\fancyhf{} % 清除页眉页脚
\cfoot{\arabic{page}} % 页脚序号
\renewcommand{\headrulewidth}{0pt} % 清除页眉页脚线
\setlength{\headheight}{15pt}

% 标题格式
\usepackage{titlesec}
\renewcommand{\thesection}{\chinese{section}}
\renewcommand{\thesubsection}{\arabic{section}.\arabic{subsection}}
\renewcommand{\thesubsubsection}{\arabic{section}.\arabic{subsection}.\arabic{subsubsection}}
\titleformat{\section}{\zihao{3}\heiti}{\thesection、}{0em}{}
\titleformat{\subsection}{\zihao{4}\heiti}{\thesubsection}{0.5em}{}
\titleformat{\subsubsection}{\zihao{4}\heiti}{\thesubsubsection}{0.5em}{}

% 代码环境
\usepackage{xcolor}
\usepackage{listings}
\usepackage{listingsutf8}
\usepackage{upquote}   % 直引号显示
\usepackage{subcaption}
\usepackage{booktabs}


\lstset{ 
  backgroundcolor=\color{gray!5},      % 浅灰背景
  frame=single,                        % 边框
  rulecolor=\color{gray!50},           % 边框颜色
  framerule=0.4pt,                     % 边框粗细
  framesep=6pt,                        % 边框与代码间距
  xleftmargin=0.5em, xrightmargin=0.5em, % 左右缩进
  aboveskip=0.8em, belowskip=0.8em,    % 上下间距
  basicstyle=\ttfamily\small,          % 等宽字体，小号
  numbers=left,                        % 显示行号
  numberstyle=\scriptsize\color{gray}, % 行号样式（更小、更淡）
  numbersep=8pt,                       % 行号和代码的间距
  breaklines=true,                     % 自动换行
  breakatwhitespace=false,             % 可在单词中间断行
  breakindent=0pt,                     % 换行后不缩进
  prebreak=\raisebox{0ex}[0ex][0ex]{\(\hookleftarrow\)\,}, % 换行前标记
  postbreak=\mbox{\space\(\hookrightarrow\)\space},        % 换行后标记
  tabsize=4,                           % Tab=4空格
  keepspaces=true,                     % 保留空格对齐
  showstringspaces=false,              % 不标记字符串里的空格
  keywordstyle=\bfseries\color{blue!70!black}, % 关键字 蓝+加粗
  commentstyle=\itshape\color{green!40!black}, % 注释 斜体墨绿
  stringstyle=\color{orange!90!black},         % 字符串 橙色
  literate=
    *{0}{{{\color{purple}0}}}{1}
     {1}{{{\color{purple}1}}}{1}
     {2}{{{\color{purple}2}}}{1}
     {3}{{{\color{purple}3}}}{1}
     {4}{{{\color{purple}4}}}{1}
     {5}{{{\color{purple}5}}}{1}
     {6}{{{\color{purple}6}}}{1}
     {7}{{{\color{purple}7}}}{1}
     {8}{{{\color{purple}8}}}{1}
     {9}{{{\color{purple}9}}}{1}
}

% 其他包
\usepackage[super, sort, compress]{cite}
\usepackage{graphicx, booktabs, tabularx}
\usepackage{amsmath}
\usepackage{amsfonts}
\usepackage{amssymb}
\usepackage{appendix}
\usepackage{caption}
\usepackage{longtable}
\usepackage{float}

\captionsetup{font=small}
\usepackage{enumitem}
\setlist[itemize]{itemsep=0.5ex, parsep=0ex, topsep=0.5ex}
\usepackage{array}
\newcolumntype{C}{>{\centering\arraybackslash}X}

% 信息
\title{}
\author{}
\date{}


%%%%%%%%%%%%%%% 正文 %%%%%%%%%%%%%%%
\begin{document}

\begin{center}
{\zihao{-2}\heiti 量化交易作业2}
\end{center}

\section{等权组合分析}
\subsection{投资策略与回测方案}

\begin{itemize}
\item \textbf{投资标的}\ 5只股票：工商银行、中国石油、万华化学、保利发展、贵州茅台
\item \textbf{初始资金}\ 10万元
\item \textbf{具体策略}\ 考虑两种买入并持有的策略：（1）将10万元全部投资于单一股票，持有至最终时刻；
（2）等权重策略：将2万元分别投资于5只股票，持有至最终时刻
\item \textbf{调仓频率}\ 买入并持有股票，期间不调仓
\item \textbf{回测区间}\ 2010年初至2021年末，共12年
\end{itemize}

\subsection{数据}

\begin{itemize}
\item 工商银行、中国石油、万华化学、保利发展、贵州茅台的收盘价，日频数据
\item 沪深300指数收盘价，日频数据
\item 各期限的国债收益率收盘价，日频数据（从财政部官网获取）
\end{itemize}

\subsection{PnL曲线}

\begin{figure}[H]
\centering
\includegraphics[width=\textwidth]{Pictures/Hw2_PnL_equal_weight.png}
\caption{等权重策略的PnL曲线（初始资金10万元）}
\end{figure}

通过PnL曲线可以看出，除中国石油外，其他4只股票和等权重策略均获得了超出市场指数（沪深300指数）的收益，
并且在2021年，多数股票出现大幅上涨时，等权重策略也表现出了类似的趋势，说明等权重组合的趋势与市场和所选股票的平均趋势具有一致性。

\subsection{策略评价指标}

首先计算各年的平均收益率和年化标准差。具体而言，先根据日频连续复利收益率计算年度收益率，
即对全年日收益率取几何平均得到年度收益率；之后计算日频连续复利收益率的标准差，
并乘以$\sqrt{252}$得到年化标准差。根据上述方法汇总得到各资产及等权重组合的年度收益率和波动率，
结果如下两表所示。


\begin{table}[H]
\centering
\caption{各年收益率}
\resizebox{\textwidth}{!}{
\begin{tabular}{lccccccc}
\toprule
年份 & 工商银行 & 中国石油 & 万华化学 & 保利发展 & 贵州茅台 & 沪深300指数 & 等权重组合 \\
\midrule
2010 & -16.14\% & -15.35\% & -17.43\% & -24.17\% &  9.24\% & -11.51\% & -12.77\% \\
2011 &   4.22\% & -10.28\% & -11.13\% &   3.98\% & 16.87\% & -25.01\% &   1.62\% \\
2012 &   2.77\% &  -3.95\% &  26.28\% &  65.62\% &  9.83\% &   7.55\% &  18.73\% \\
2013 &  -8.17\% & -11.67\% &  37.72\% & -38.20\% & -36.56\% &  -7.65\% & -15.59\% \\
2014 &  46.29\% &  46.30\% &   9.83\% & 104.74\% & 66.94\% &  51.66\% &  50.57\% \\
2015 &  -1.53\% & -21.58\% & -17.18\% &  -0.04\% & 28.81\% &   5.58\% &  -0.53\% \\
2016 &   1.37\% &  -4.27\% &  22.07\% & -10.81\% & 56.45\% & -11.28\% &  17.17\% \\
2017 &  47.06\% &   3.19\% & 115.28\% &  59.92\% &111.89\% &  21.78\% &  82.80\% \\
2018 & -10.90\% &  -9.24\% & -23.19\% & -14.19\% &-14.21\% & -25.31\% & -15.46\% \\
2019 &  16.05\% & -17.04\% & 111.20\% &  42.87\% &103.47\% &  36.07\% &  76.72\% \\
2020 & -10.71\% & -26.06\% &  66.79\% &   3.24\% & 70.86\% &  27.21\% &  51.22\% \\
2021 &  -2.13\% &  23.03\% &  12.25\% &   5.19\% &  3.57\% &  -5.20\% &   5.90\% \\
\bottomrule
\end{tabular}
}
\end{table}

从各年平均收益率表可以看出，在收益率表现方面，等权重组合在多数年份能够获得较为可观的收益，
并且在市场上涨阶段通常能够取得高于市场指数的表现。例如在2012、2014、2017、2019和2020年，
等权重组合的收益率均明显高于沪深300指数，这表明组合能够较好地捕捉股票市场整体上涨所带来的收益。

当市场出现下跌时，等权重策略在部分年份也表现出一定的抗跌能力。例如在2011、2016和2021年，
等权重组合的收益率均优于沪深300指数。此外，在2018年市场整体大幅下跌的情况下，
等权重组合的跌幅仍小于市场指数。总体来看，组合收益虽然在个别年份仍受到股票池个股表现的影响，
但整体表现优于市场指数，说明在所选股票池条件下等权重配置具有一定的收益优势。

从长期趋势来看，等权重组合在2017年之后的收益表现尤为突出，
例如2017年和2019年的组合收益率分别达到82.80\%和76.72\%。
这一现象主要受到股票池中部分成长性较强股票（如万华化学和贵州茅台）表现优异的带动，
从而显著提升了组合整体收益水平。

\begin{table}[H]
\centering
\caption{各年日收益率标准差}
\resizebox{\textwidth}{!}{
\begin{tabular}{lccccccc}
\toprule
年份 & 工商银行 & 中国石油 & 万华化学 & 保利发展 & 贵州茅台 & 沪深300指数 & 等权重组合 \\
\midrule
2010 & 1.52\% & 1.52\% & 2.71\% & 2.89\% & 1.94\% & 1.59\% & 1.46\% \\
2011 & 1.03\% & 0.99\% & 2.47\% & 2.43\% & 1.55\% & 1.30\% & 1.19\% \\
2012 & 0.83\% & 0.80\% & 1.88\% & 2.36\% & 1.96\% & 1.29\% & 1.18\% \\
2013 & 1.00\% & 0.95\% & 2.06\% & 2.44\% & 1.86\% & 1.40\% & 1.21\% \\
2014 & 1.50\% & 1.50\% & 1.84\% & 2.57\% & 1.96\% & 1.21\% & 1.28\% \\
2015 & 2.45\% & 3.07\% & 3.40\% & 3.94\% & 2.52\% & 2.48\% & 2.51\% \\
2016 & 0.87\% & 1.24\% & 2.25\% & 2.32\% & 1.67\% & 1.40\% & 1.37\% \\
2017 & 1.05\% & 0.86\% & 2.38\% & 1.77\% & 1.70\% & 0.64\% & 1.00\% \\
2018 & 1.62\% & 1.64\% & 2.54\% & 2.86\% & 2.24\% & 1.35\% & 1.72\% \\
2019 & 0.98\% & 0.84\% & 2.21\% & 1.95\% & 1.91\% & 1.25\% & 1.49\% \\
2020 & 1.12\% & 1.25\% & 2.58\% & 2.03\% & 1.83\% & 1.43\% & 1.56\% \\
2021 & 0.82\% & 1.82\% & 2.78\% & 2.52\% & 2.36\% & 1.17\% & 1.81\% \\
\bottomrule
\end{tabular}
}
\end{table}

而根据各年收益率标准差表可以看出，虽然所选择的5只股票中存在波动率较高的股票，
例如万华化学和保利发展，但等权重组合的整体波动率并未明显高于市场指数，
甚至在部分年份低于市场指数，例如2010、2011、2012和2016年。
这一现象体现了投资组合分散化的作用，即通过同时配置多只股票，
可以在一定程度上降低单只股票价格波动带来的风险，从而使组合整体波动更加稳定。

随后，我们进一步计算整个回测区间内的综合绩效指标，包括年化收益率、年化标准差、
夏普比率以及最大回撤等指标，结果如下表所示。

\begin{table}[H]
\centering
\caption{策略评价指标（1）}
\begin{tabular}{lcccc}
\toprule
资产 & 年化收益率 & 年化标准差 & 夏普比率 & 最大回撤（万元） \\
\midrule
工商银行 & 4.15\% & 20.82\% & 0.0674 & 7.63 \\
中国石油 & -5.81\% & 23.84\% & -0.3587 & 8.81 \\
万华化学 & 21.22\% & 39.05\% & 0.4731 & 43.39 \\
保利发展 & 10.38\% & 40.69\% & 0.1878 & 15.34 \\
贵州茅台 & 29.32\% & 31.43\% & 0.8457 & 98.30 \\
沪深300指数 & 2.93\% & 22.75\% & 0.0085 & 7.07 \\
等权重策略 & 18.04\% & 24.30\% & 0.6298 & 26.40 \\
\bottomrule
\end{tabular}
\end{table}

从表中可以看出，等权重策略的年化收益率为18.04\%，明显高于沪深300指数的2.93\%。
与此同时，两者的年化波动率分别为24.30\%和22.75\%，差异并不显著。
因此在风险调整收益方面，等权重策略的夏普比率达到0.6298，
显著高于市场指数的0.0085，说明在单位风险下该策略能够获得更高的收益，
在均值—方差框架下具有更优的风险收益特征。

但从最大回撤来看，等权重策略的最大回撤达到26.40万元，
明显高于市场指数的7.07万元，且初始资金只有10万元。
这表明在市场出现较大波动时，组合仍然可能面临较大的阶段性亏损，
因此在风险控制方面仍存在一定改进空间。

\begin{table}[H]
\centering
\caption{策略评价指标（2）}
\begin{tabular}{lcccc}
\toprule
资产 & $\beta$ & 特雷诺比率 & 詹森$\alpha$ & 信息比率 \\
\midrule
工商银行 & 0.5079 & 0.0277 & 0.0131 & 0.0754 \\
中国石油 & 0.5967 & -0.1433 & -0.0867 & -0.4421 \\
万华化学 & 1.1064 & 0.1670 & 0.1826 & 0.6117 \\
保利发展 & 1.1828 & 0.0646 & 0.0741 & 0.2429 \\
贵州茅台 & 0.7307 & 0.3637 & 0.2643 & 0.9912 \\
沪深300指数 & 1.0000 & 0.0019 & 0.0000 & 0.0088 \\
等权重组合 & 0.8824 & 0.1734 & 0.1513 & 1.1054 \\
\bottomrule
\end{tabular}
\end{table}

继续看如上第二张表，从$\beta$值来看，等权重组合的$\beta$为0.8824，略低于市场指数的1，
说明该组合对市场整体波动的敏感程度略低于市场，
在一定程度上具有降低系统性风险的作用。

从特雷诺比率来看，等权重组合的特雷诺比率为0.1734，
高于除贵州茅台之外的其他股票，
说明在考虑系统性风险后，该策略能够获得较为可观的超额收益。

在信息比率方面，等权重组合的信息比率达到1.1054，
明显高于大多数单只股票，
说明相对于市场指数，该组合能够获得更加稳定的超额收益。

此外，从詹森$\alpha$指标来看，除中国石油外，
其余股票以及等权重组合的$\alpha$均为正值，
说明在考虑市场风险之后仍然能够获得正的超额收益，
其中贵州茅台和等权重组合表现尤为突出。

综合收益与风险指标可以看出，
等权重组合在长期收益率、风险调整收益以及相对收益稳定性方面均明显优于市场指数，
说明在所选股票池条件下，
简单的等权配置策略已经能够取得较好的投资效果。


\section{基于认知规则和基于模型的信号}

\subsection{投资策略与回测方案}

\begin{itemize}
\item \textbf{投资标的，初始资金}\ 与第一问相同
\item \textbf{具体策略}\ 考虑两种策略：
（1）\textbf{策略1（认知）}：将资金等分为5份，每只股票分配2万元。
对每只股票分别计算其20日和60日移动平均线，
若20日均线高于60日均线，则认为该股票处于上升趋势，
在下一交易日持有该股票；否则不持有该股票。
组合根据各股票的持仓信号进行投资，并在每日根据新的信号进行调整。
（2）\textbf{策略2（模型）}：对每只股票分别建立随机森林分类模型。
具体地，以该股票过去20个交易日的收益率序列作为输入特征，
以该股票下一交易日的涨跌方向作为预测目标：若下一交易日收益率大于0，则记为1，否则记为0。
样本划分为训练集、验证集和测试集，其中训练集与验证集用于模型训练和参数选择，
测试集用于样本外回测。
在模型训练过程中，采用随机森林分类器，并结合时间序列交叉验证（TimeSeriesSplit）
和网格搜索（GridSearchCV）对超参数进行调优，
以F1值作为评价标准。
在测试阶段，对5只股票分别生成每日涨跌预测信号，
并根据预测结果构造交易策略，每日调整持仓。
\item \textbf{调仓频率}\ 日频
\item \textbf{回测区间}\ 2017年3月22日至2021年末（测试集占全样本的40\%）
\end{itemize}

\subsection{数据}

\begin{itemize}
\item 工商银行、中国石油、万华化学、保利发展、贵州茅台的收盘价，日频数据
\item 沪深300指数收盘价，日频数据
\item 各期限的国债收益率收盘价，日频数据（从财政部官网获取）
\end{itemize}

\subsection{PnL曲线}

\begin{figure}[H]
\centering
\includegraphics[width=\textwidth]{Pictures/Hw2_PnL_1_2.png}
\caption{两种策略的PnL曲线（初始资金10万元）}
\end{figure}

从PnL曲线可以看出，在整个样本区间内，策略1与策略2的资金曲线整体均呈现稳步上升趋势，且长期表现均优于沪深300指数。与指数相比，两种策略的净值波动明显更小，资金曲线也更加平滑，说明基于交易信号进行择时能够在一定程度上降低组合对市场整体波动的暴露，从而改善策略的风险收益特征。

从两种策略之间的比较来看，策略1在样本期内的累计收益略高于策略2，期末资金规模也更高；但与此同时，策略2的资金曲线整体更为平稳，回撤控制相对更好。特别是在2018年市场整体处于下行阶段时，两种策略都明显跑赢沪深300指数，其中策略2跌幅更小，表现出更强的防御性；而在2019年至2020年的上涨阶段，策略1的净值增长速度更快，体现出其在趋势行情中更强的收益弹性。

同时，从图中还可以看出，等权重组合在样本区间内取得了最高的绝对收益，其资金曲线在2019年至2021年期间上升最为明显，一度显著高于两种择时策略和沪深300指数。但与此同时，等权重组合的波动也更大，在2021年前后出现了较为明显的回撤。这表明，直接持有股票池虽然能够更充分地分享上涨行情带来的收益，但也更容易受到组合中高波动个股的影响。

相比之下，策略1和策略2的优势主要体现在收益的稳健性和风险控制能力上，而不是绝对收益水平最高。两种策略的资金曲线在大部分时间内均高于沪深300指数，并且波动幅度显著低于等权重组合，说明无论是基于均线规则构建的择时信号，还是基于随机森林模型生成的交易信号，都能够在一定程度上提升资金使用效率，减少无效持仓和市场下跌阶段带来的损失。


\subsection{策略评价指标}

同第一问，首先计算两种策略各年的平均收益率和收益率标准差，并与沪深300指数和等权重组合进行对比，结果如下。

\begin{table}[H]
\centering
\caption{各年收益率}
\begin{tabular}{lcccc}
\toprule
年份 & 沪深300指数 & 等权重组合 & 策略1 & 策略2 \\
\midrule
2017 & 17.29\% & 62.79\% & 39.30\% & 18.37\% \\
2018 & -25.31\% & -15.46\% & -5.32\% & -2.88\% \\
2019 & 36.07\% & 76.72\% & 26.52\% & 18.33\% \\
2020 & 27.21\% & 51.22\% & 11.34\% & 8.93\% \\
2021 & -5.20\% & 5.90\% & 5.91\% & 15.21\% \\
\bottomrule
\end{tabular}
\end{table}

从各年平均收益率可以看出，两种策略在多数年份均取得了正收益，且整体表现明显优于沪深300指数。具体来看，2017年策略1和策略2的收益率分别为39.30\%和18.37\%，均高于沪深300指数的17.29\%；2018年在市场大幅下跌的背景下，沪深300指数收益率为-25.31\%，等权重组合也为-15.46\%，而策略1和策略2分别仅为-5.32\%和-2.88\%，跌幅明显更小，体现出较强的抗跌能力；2019年和2020年市场上涨阶段，两种策略继续保持正收益，虽然收益水平低于等权重组合，但整体表现依然较为稳健；到了2021年，在沪深300指数收益率为-5.20\%的情况下，策略1和策略2仍分别取得5.91\%和15.21\%的正收益，显示出信号策略在震荡股市中仍具备较好的适应性。

整体来看，等权重组合在上涨年份往往能够取得更高的收益水平，但其收益波动也更为明显；相比之下，策略1和策略2虽然绝对收益并非始终最高，但年度收益表现更加均衡。

\begin{table}[H]
\centering
\caption{各年日收益率标准差}
\begin{tabular}{lcccc}
\toprule
年份 & 沪深300指数 & 等权重组合 & 策略1 & 策略2 \\
\midrule
2017 & 0.67\% & 1.06\% & 0.71\% & 0.53\% \\
2018 & 1.35\% & 1.72\% & 0.92\% & 0.88\% \\
2019 & 1.25\% & 1.49\% & 1.00\% & 0.61\% \\
2020 & 1.43\% & 1.56\% & 0.96\% & 0.75\% \\
2021 & 1.17\% & 1.81\% & 0.84\% & 0.78\% \\
\bottomrule
\end{tabular}
\end{table}

根据各年收益率标准差可以看出，两种策略的波动率整体显著低于等权重组合，并且在绝大多数年份也低于沪深300指数。具体而言，2018年至2021年期间，策略1和策略2的收益率标准差始终低于市场指数；其中策略2在各年份中的波动率均处于较低水平，例如2019年仅为0.61\%，明显低于策略1的1.00\%、沪深300指数的1.25\%以及等权重组合的1.49\%。这说明通过交易信号进行择时，能够有效降低组合暴露于市场波动的程度，从而改善策略的稳定性。

总体来看，策略1和策略2在收益稳定性方面均优于等权重组合，其中策略2的年度波动控制更为突出，表现出更明显的低波动特征。

\begin{table}[H]
\centering
\caption{策略评价指标}
\begin{tabular}{lcccc}
\toprule
指标 & 沪深300指数 & 等权重组合 & 策略1 & 策略2 \\
\midrule
平均年化收益率 & 8.20\% & 34.36\% & 15.84\% & 12.32\% \\
年化标准差 & 19.48\% & 24.92\% & 14.31\% & 11.59\% \\
夏普比率 & 0.2804 & 1.2688 & 0.9156 & 0.8269 \\
$\beta$ & 1.0000 & 1.0194 & 0.5217 & 0.3891 \\
特雷诺比率 & 0.0546 & 0.3102 & 0.2511 & 0.2463 \\
詹森$\alpha$ & 0.0000 & 0.2605 & 0.1025 & 0.0746 \\
信息比率 & 0.3024 & 1.7296 & 1.0176 & 0.8505 \\
最大回撤（万元） & 4.15 & 15.08 & 3.06 & 1.29 \\
\bottomrule
\end{tabular}
\end{table}

最后进一步比较整个回测区间内的综合绩效指标。从表中可以看出，策略1和策略2的平均年化收益率分别为15.84\%和12.32\%，均明显高于沪深300指数的8.20\%，说明两种策略在长期表现上均能够获得超越市场的收益。但与等权重组合34.36\%的年化收益相比，两种择时策略在绝对收益上仍存在一定差距，这与PnL曲线中等权重组合在牛市阶段上涨更快的现象是一致的。

从风险水平来看，策略1和策略2的年化标准差分别为14.31\%和11.59\%，显著低于市场指数的19.48\%以及等权重组合的24.92\%。这说明两种策略在控制波动方面具有明显优势，尤其是策略2，其整体波动水平最低。因此，从风险调整后的收益角度看，两种策略都优于市场指数，但策略1的夏普比率为0.9156，高于策略2的0.8269，说明虽然策略2更稳健，但策略1凭借更高的年化收益，在单位总风险上的收益补偿更高。

在市场风险暴露方面，策略1和策略2的$\beta$分别为0.5217和0.3891，均显著低于1，说明两种策略对市场整体波动的敏感度较低，具有较好的防御属性。对应地，两种策略的特雷诺比率分别为0.2511和0.2463，均显著高于沪深300指数的0.0546，表明在承担系统性风险的情况下，两种策略均实现了更高的超额回报。从数值上看，策略1的特雷诺比率略高于策略2，说明其在单位系统性风险上的收益表现略优。

从詹森$\alpha$来看，策略1和策略2分别为0.1025和0.0746，均为正值，说明在考虑市场风险之后，两种策略仍然能够获得正的超额收益。不过，策略1的$\alpha$高于策略2，表明其相对于CAPM基准所体现出的主动收益更强。

在信息比率方面，策略1和策略2分别为1.0176和0.8505，均显著高于沪深300指数的0.3024，说明两种策略相对于基准指数均具有较好的超额收益稳定性。其中策略1的信息比率更高，表明其相对基准获得超额收益的效率更优。

最后，从最大回撤来看，策略1和策略2的最大回撤分别为3.06万元和1.29万元，均明显低于等权重组合的15.08万元，也低于沪深300指数的4.15万元。其中策略2的最大回撤最小，说明其在控制极端亏损方面具有最强的优势，这一点也与PnL曲线中其资金轨迹更加平滑的特征相一致。

综合以上分析可以看出，两种基于信号的交易策略在收益与风险的综合表现上均优于市场指数。若从绝对收益、夏普比率、特雷诺比率、詹森$\alpha$和信息比率等指标来看，策略1整体表现略优；但若从波动率、$\beta$以及最大回撤等风险控制指标来看，策略2更具优势。因此，策略1更适合追求较高风险调整收益的投资者，而策略2则更适合偏好稳健、重视回撤控制的投资者。

\end{document}